% GENERATED FILE. Edit canonical lessons, then run npm run generate:books.
% canonical-source-hash: fnv1a64:b8b5e9eefa45db07
% canonical-lessons: LA-C476-trituro, LA-C476-ventilo, LA-C476-palea, LA-C476-stipula, LA-C476-horreum

\chapter{To thresh, To winnow, Chaff, Stalk, Granary}
\label{ch:la-to-thresh-to-winnow-chaff-stalk-granary}
\hlchaptermodality{476}

\begin{chapteropening}
\textbf{By the end of this chapter:} \emph{I can talk about what happens after the harvest: threshing, winnowing, the chaff, the stalks and storing grain in the granary.}

Complete the last lesson of chapter 476: I can talk about what happens after the harvest: threshing, winnowing, the chaff, the stalks and storing grain in the granary.
\end{chapteropening}

\section[trītūrō, trītūrāre]{trītūrō, trītūrāre — to thresh}

\label{lesson:LA-C476-trituro}



\begin{quote}
Before the new one: say the Latin for an ear of grain, then the Latin for wheat.
\end{quote}

\subsection*{You'll want to know: trītūrō, trītūrāre}
\textbf{trītūrō, trītūrāre} — \textquotedblleft{}to thresh\textquotedblright{}.

\begin{quote}
In āreā servī frūmentum trītūrant.
\end{quote}

\begin{quote}
\emph{On the threshing floor the slaves thresh the grain.}
\end{quote}

\emph{trītūrō, trītūrāre, trītūrāvī, trītūrātum}, from \emph{trītūra} 'threshing'. Oxen often trod the grain out. A late word; earlier writers say \emph{frūmentum terere}.

\begin{cousinweb}
\emph{trītūrō} is built on \emph{trītūra} 'threshing' (EtymDB; Lewis \& Short), which comes from \emph{terō} 'to rub'. English \emph{triturate} 'to grind to powder' comes from it.
\end{cousinweb}

\begin{grammarlens}[title={What you've built}]
One more word for telling, explaining and describing at B1.
\end{grammarlens}

\subsection*{Your turn}
\begin{itemize}
\raggedright
  \item \emph{Say it:} \emph{trītūrō, trītūrāre}
  \item \emph{Say it:} \emph{trītūrō, trītūrāre}, once more
  \item \emph{Say it:} a sentence of your own with \emph{trītūrō, trītūrāre}
  \item \emph{Recall:} say the Latin for an ear of grain, then the Latin for wheat, then say \emph{trītūrō, trītūrāre} again
\end{itemize}

\begin{tcolorbox}[breakable,colback=teal!4,colframe=teal!35!black,title={Before you move on}]
What does trītūrō, trītūrāre mean? (\textquotedblleft{}To thresh\textquotedblright{}.) Say it once more.
\end{tcolorbox}
\section[ventilō, ventilāre]{ventilō, ventilāre — to winnow; to fan, to toss in the air}

\label{lesson:LA-C476-ventilo}



\begin{quote}
Before the new one: say the Latin for wheat, then the Latin for to thresh.
\end{quote}

\subsection*{You'll want to know: ventilō, ventilāre}
\textbf{ventilō, ventilāre} — \textquotedblleft{}to winnow; to fan, to toss in the air\textquotedblright{}.

\begin{quote}
Agricolae frūmentum in āreā ventilant.
\end{quote}

\begin{quote}
\emph{The farmers winnow the grain on the threshing floor.}
\end{quote}

\emph{ventilō, ventilāre, ventilāvī, ventilātum}. English \emph{ventilate} comes from it.

\begin{cousinweb}
\emph{ventilō} is built on \emph{ventulus} 'a slight wind' (EtymDB; Lewis \& Short), a diminutive of \emph{ventus} 'wind'. English \emph{ventilate} comes from it.
\end{cousinweb}

\begin{grammarlens}[title={What you've built}]
One more word for telling, explaining and describing at B1.
\end{grammarlens}

\subsection*{Your turn}
\begin{itemize}
\raggedright
  \item \emph{Say it:} \emph{ventilō, ventilāre}
  \item \emph{Say it:} \emph{ventilō, ventilāre}, once more
  \item \emph{Say it:} a sentence of your own with \emph{ventilō, ventilāre}
  \item \emph{Recall:} say the Latin for wheat, then the Latin for to thresh, then say \emph{ventilō, ventilāre} again
\end{itemize}

\begin{tcolorbox}[breakable,colback=teal!4,colframe=teal!35!black,title={Before you move on}]
What does ventilō, ventilāre mean? (\textquotedblleft{}To winnow; to fan, to toss in the air\textquotedblright{}.) Say it once more.
\end{tcolorbox}
\section[palea, paleae]{palea, paleae — chaff; straw}

\label{lesson:LA-C476-palea}



\begin{quote}
Before the new one: say the Latin for to thresh, then the Latin for to winnow.
\end{quote}

\subsection*{You'll want to know: palea, paleae}
\textbf{palea, paleae} — \textquotedblleft{}chaff; straw\textquotedblright{}.

\begin{quote}
Palea levis est et ā ventō portātur.
\end{quote}

\begin{quote}
\emph{Chaff is light and is carried by the wind.}
\end{quote}

First declension, feminine: \emph{palea, paleae f.} The husks left after threshing, used for fodder and bedding. English keeps \emph{palea} as a term in botany.

\begin{cousinweb}
\emph{palea} is inherited from PIE \emph{pel-} 'flour, dust' (EtymDB), the root also seen in \emph{pulvis} 'dust' and \emph{pollen}. French \emph{paille} 'straw' comes from it.
\end{cousinweb}

\begin{grammarlens}[title={What you've built}]
One more word for telling, explaining and describing at B1.
\end{grammarlens}

\subsection*{Your turn}
\begin{itemize}
\raggedright
  \item \emph{Say it:} \emph{palea, paleae}
  \item \emph{Say it:} \emph{palea, paleae}, once more
  \item \emph{Say it:} a sentence of your own with \emph{palea, paleae}
  \item \emph{Recall:} say the Latin for to thresh, then the Latin for to winnow, then say \emph{palea, paleae} again
\end{itemize}

\begin{tcolorbox}[breakable,colback=teal!4,colframe=teal!35!black,title={Before you move on}]
What does palea, paleae mean? (\textquotedblleft{}Chaff; straw\textquotedblright{}.) Say it once more.
\end{tcolorbox}
\section[stipula, stipulae]{stipula, stipulae — a stalk, a straw; stubble}

\label{lesson:LA-C476-stipula}



\begin{quote}
Before the new one: say the Latin for to winnow, then the Latin for chaff.
\end{quote}

\subsection*{You'll want to know: stipula, stipulae}
\textbf{stipula, stipulae} — \textquotedblleft{}a stalk, a straw; stubble\textquotedblright{}.

\begin{quote}
Post messem in agrō sōla stipula manet.
\end{quote}

\begin{quote}
\emph{After the harvest only stubble is left in the field.}
\end{quote}

First declension, feminine: \emph{stipula, stipulae f.}, a diminutive linked with \emph{stīpes} 'stake'.

\begin{cousinweb}
\emph{stipula} is inherited from Proto-Italic \emph{stipelā} 'straw' (EtymDB), and Lewis \& Short connect it with \emph{stīpes} 'stake, trunk'. English \emph{stubble} comes, through French, from a related Late Latin form.
\end{cousinweb}

\begin{grammarlens}[title={What you've built}]
One more word for telling, explaining and describing at B1.
\end{grammarlens}

\subsection*{Your turn}
\begin{itemize}
\raggedright
  \item \emph{Say it:} \emph{stipula, stipulae}
  \item \emph{Say it:} \emph{stipula, stipulae}, once more
  \item \emph{Say it:} a sentence of your own with \emph{stipula, stipulae}
  \item \emph{Recall:} say the Latin for to winnow, then the Latin for chaff, then say \emph{stipula, stipulae} again
\end{itemize}

\begin{tcolorbox}[breakable,colback=teal!4,colframe=teal!35!black,title={Before you move on}]
What does stipula, stipulae mean? (\textquotedblleft{}A stalk, a straw; stubble\textquotedblright{}.) Say it once more.
\end{tcolorbox}
\section[horreum, horreī]{horreum, horreī — a granary, a barn, a storehouse}

\label{lesson:LA-C476-horreum}



\begin{quote}
Before the new one: say the Latin for chaff, then the Latin for a stalk.
\end{quote}

\subsection*{You'll want to know: horreum, horreī}
\textbf{horreum, horreī} — \textquotedblleft{}a granary, a barn, a storehouse\textquotedblright{}.

\begin{quote}
Frūmentum in horreō servātur.
\end{quote}

\begin{quote}
\emph{The grain is kept in the granary.}
\end{quote}

Second declension, neuter: \emph{horreum, horreī n.} Big public \emph{horrea} by the river stored the city's grain.

\begin{cousinweb}
The origin of \emph{horreum} is not settled. EtymDB traces it to Greek \emph{hōreion} 'storehouse, granary', while Lewis \& Short compare \emph{farreum}, from \emph{far} 'spelt'.
\end{cousinweb}

\begin{grammarlens}[title={What you've built}]
One more word for telling, explaining and describing at B1.
\end{grammarlens}

\subsection*{Your turn}
\begin{itemize}
\raggedright
  \item \emph{Say it:} \emph{horreum, horreī}
  \item \emph{Say it:} \emph{horreum, horreī}, once more
  \item \emph{Say it:} a sentence of your own with \emph{horreum, horreī}
  \item \emph{Recall:} say the Latin for chaff, then the Latin for a stalk, then say \emph{horreum, horreī} again
\end{itemize}

\begin{tcolorbox}[breakable,colback=teal!4,colframe=teal!35!black,title={Before you move on}]
What does horreum, horreī mean? (\textquotedblleft{}A granary, a barn, a storehouse\textquotedblright{}.) Say it once more.
\end{tcolorbox}
